\documentclass[a4paper,oneside,11pt]{book}

% Package yang diperlukan
\usepackage{graphicx}
\usepackage{titling}
\usepackage{geometry}
\usepackage{listings}
\usepackage{xcolor}
\usepackage{float}
\usepackage{hyperref}
\usepackage{indentfirst}
\usepackage{array}
\usepackage[table]{xcolor}
\usepackage{tabularx}
\usepackage[indonesian]{babel}
\usepackage{multirow}
\usepackage{booktabs}
\usepackage{url}
\usepackage{mathptmx}   
\usepackage{helvet}    
\usepackage{amsfonts}   
\usepackage{tikz}
\usetikzlibrary{shapes.geometric, arrows}

% Untuk manajemen font dan karakter khusus
\usepackage[T1]{fontenc}       % Encoding font Latin modern
\usepackage{textcomp}          % Simbol tambahan (text companion)

% Untuk referensi silang yang lebih baik
\usepackage[all]{hypcap}       % Perbaiki link ke gambar/tabel

% Untuk kontrol paragraf
\usepackage{parskip}           % Spasi antar paragraf konsisten

% Untuk pengaturan halaman
\usepackage{fancyhdr}          % Header/footer custom
\pagestyle{fancy}              % Aktifkan style fancy
\fancyhf{}                     % Kosongkan header/footer
\rfoot{\thepage}               % Nomor halaman di kanan bawah

% Untuk gambar yang lebih fleksibel
\usepackage{adjustbox}         % Penyesuaian ukuran gambar otomatis

% Untuk kode program tambahan
\usepackage{etoolbox}          % Tools untuk programming LaTeX

\hypersetup{
    colorlinks=true,
    linkcolor=black,
    citecolor=blue,
    urlcolor=blue
}

% Pengaturan margin
\geometry{
    left=3cm,
    right=2.5cm,
    top=3cm,
    bottom=3cm
}

% Pengaturan untuk kode HTML/Python
\definecolor{codegreen}{rgb}{0,0.6,0}
\definecolor{codegray}{rgb}{0.5,0.5,0.5}
\definecolor{codepurple}{rgb}{0.58,0,0.82}
\definecolor{backcolour}{rgb}{0.95,0.95,0.92}

\lstdefinestyle{htmlstyle}{
    backgroundcolor=\color{backcolour},   
    commentstyle=\color{codegreen},
    keywordstyle=\color{codepurple},
    numberstyle=\tiny\color{codegray},
    stringstyle=\color{codegreen},
    basicstyle=\ttfamily\footnotesize,
    breakatwhitespace=false,         
    breaklines=true,                 
    captionpos=b,                    
    keepspaces=true,                 
    numbers=left,                    
    numbersep=5pt,                  
    showspaces=false,                
    showstringspaces=false,
    showtabs=false,                  
    tabsize=2,
    frame=single
}

\begin{document}

% ================= HALAMAN JUDUL =================
\begin{titlingpage}
\begin{center}
% Spasi atas
\vspace*{0.5cm}

% Bagian Judul Atas
{\large \textbf{LAPORAN PRAKTIKUM}} \\[0.3cm]
{\large \textbf{PRAKTIKUM PENAMBANGAN DATA}} \\[0.3cm]
{\large \textbf{PERTEMUAN 2}} \\[0.3cm]
{\large \textbf{Exploratory Data Analysis (EDA) dan Data Preprocessing}} \\[2.5cm] % Ubah judul materi sesuai kebutuhan

% Bagian Logo
% Pastikan file lambangugm.png berada di satu folder dengan file .tex
\includegraphics[width=6cm]{lambangugm.png} \\[2.5cm]

% Bagian Identitas
\normalsize
Disusun oleh:\\[0.5cm]
\begin{tabular}{@{}l@{\hspace{1cm}}l@{}}
Nama           & : Hilarius Christiano Avin Paliling \\
NIM            & : 24/542159/SV/25009 \\
Kelas          & : PL5B1 \\
Dosen Pengampu & : Dr.Eng. Ir. Ganjar Alfian, S.T., M.Eng.
\end{tabular}

% vfill akan secara otomatis mendorong teks di bawahnya ke dasar halaman
\vfill 

% Bagian Institusi Bawah
{\large \textbf{PROGRAM STUDI D-IV TEKNOLOGI REKAYASA PERANGKAT LUNAK}} \\[0.1cm]
{\large \textbf{DEPARTEMEN TEKNIK ELEKTRO DAN INFORMATIKA}} \\[0.1cm]
{\large \textbf{SEKOLAH VOKASI}} \\[0.1cm]
{\large \textbf{UNIVERSITAS GADJAH MADA}} \\[0.1cm]
{\large \textbf{YOGYAKARTA}} \\[0.1cm]
{\large \textbf{2026}}

% Spasi batas bawah agar tidak menyentuh margin
\vspace*{0.5cm}
\end{center}
\end{titlingpage}
% =================================================

% Daftar Isi
\tableofcontents
\listoffigures
\newpage

\chapter{Tujuan Praktikum}
\begin{enumerate}
    \item Mahasiswa mampu memahami konsep dasar Artificial Intelligence (AI)
    \item Mahasiswa mengetahui dasar-dasar Machine Learning (ML)
    \item Mahasiswa dapat menerapkan langkah awal analisis data
    \item Mahasiswa mampu mengetahui dan menerapkan teknik-teknik dasar dalam data preprocessing
\end{enumerate}

\chapter{Dasar Teori}
% ==========================================
% ARTIFICIAL INTELLIGENCE & MACHINE LEARNING
% ==========================================
\section{Artificial Intelligence dan Machine Learning}
\textit{Artificial Intelligence} (AI) merupakan cabang ilmu komputer yang bertujuan merekayasa sistem komputasi agar mampu meniru kemampuan kognitif manusia, seperti penalaran, pengenalan pola, serta pengambilan keputusan adaptif [\ref{modulPraktikum}]. Perkembangan AI modern sangat berpusat pada data (\textit{data-driven}), di mana performa sistem ditentukan secara langsung oleh kualitas data yang dihimpun dan diproses.
Sebagai sub-bidang inti dari AI, \textit{Machine Learning} (ML) menyediakan metodologi yang memungkinkan komputer belajar secara otomatis dari data historis tanpa harus diprogram secara eksplisit [\ref{mckinney}]. Pendekatan algoritma ML secara umum terbagi menjadi tiga paradigma:
\begin{enumerate}
    \item \textbf{\textit{Supervised Learning}}: Menggunakan pasangan data fitur dan label target untuk tugas prediksi atau klasifikasi.
    \item \textbf{\textit{Unsupervised Learning}}: Menemukan pola tersembunyi, klasterisasi, atau anomali pada data tanpa label.
    \item \textbf{\textit{Reinforcement Learning}}: Mengoptimalkan tindakan melalui mekanisme umpan balik berupa penghargaan (\textit{reward}) dan penalti (\textit{penalty}).
\end{enumerate}
Keberhasilan implementasi model ML berbanding lurus dengan ketepatan perlakuan data sebelum proses pembelajaran dilakukan [\ref{modulPraktikum}].
\begin{figure}[H]
    \centering
    \includegraphics[width=0.6\textwidth]{60.jpg}
    \caption{Domain Artificial Intelligence dan taksonomi pembagian Machine Learning}
    \label{fig:60}
\end{figure}
% ==========================================
% EXPLORATORY DATA ANALYSIS (EDA)
% ==========================================
\clearpage
\section{Exploratory Data Analysis}
Sebelum data diumpankan ke dalam algoritma pemodelan, tahap esensial yang wajib dilaksanakan adalah \textit{Exploratory Data Analysis} (EDA) [\ref{modulPraktikum}]. Konsep EDA yang dipelopori oleh John W. Tukey menekankan urgensi menginspeksi data secara visual dan statistik deskriptif sebelum merumuskan asumsi atau hipotesis lanjutan [\ref{tukey}].
Dalam praktiknya, alur EDA berfokus pada evaluasi struktur data melalui metrik kuartil, tendensi sentral (\textit{mean}, \textit{median}), identifikasi nilai ekstrem (\textit{outliers}), serta deteksi anomali seperti nilai kosong (\textit{missing values}) dan duplikasi data. Selain ringkasan statistik numerik, visualisasi data memegang peranan krusial melalui penggunaan berbagai diagram:
\begin{itemize}
    \item \textbf{\textit{Histogram \& Box Plot}}: Memetakan sebaran distribusi frekuensi variabel numerik serta mendeteksi pencilan data.
    \item \textbf{\textit{Bar Chart \& Pie Chart}}: Meninjau proporsi serta persebaran frekuensi pada fitur kategorikal.
    \item \textbf{\textit{Scatter Plot \& Heatmap Korelasi}}: Mengukur keeratan asosiasi atau derajat korelasi Pearson antar variabel.
\end{itemize}
Tahapan ini menjadi landasan diagnosa awal untuk menentukan strategi pembersihan data yang tepat [\ref{mckinney}].
\begin{figure}[H]
    \centering
    \includegraphics[width=0.8\textwidth]{61.png}
    \caption{Visualisasi analitik korelasi dan distribusi data dalam proses EDA}
    \label{fig:61}
\end{figure}
% ==========================================
% DATA PREPROCESSING
% ==========================================
\clearpage
\section{Data Preprocessing}
\textit{Data preprocessing} adalah proses rekayasa sistematis untuk membersihkan, mentransformasikan, dan menyusun data mentah menjadi representasi matematis yang siap diolah oleh algoritma [\ref{modulPraktikum}]. Kualitas prapemrosesan ini menentukan ketahanan model terhadap bias dan \textit{overfitting}.
Alur kerja \textit{data preprocessing} mencakup beberapa tahapan utama:
\begin{enumerate}
    \item \textbf{\textit{Data Cleaning}}: Meliputi eliminasi duplikasi baris serta penanganan \textit{missing values} melalui metode penghapusan (\textit{dropping}) atau imputasi nilai pengganti (misalnya \textit{median} untuk variabel numerik atau modus untuk variabel kategorikal).
    \item \textbf{\textit{Categorical Encoding}}: Mengubah data string/teks ke format numerik agar dapat diproses oleh algoritma komputasi. Teknik yang umum digunakan meliputi \textit{Label Encoding} (untuk data ordinal/sederhana) dan \textit{One-Hot Encoding} (untuk data nominal tanpa tingkatan).
    \item \textbf{\textit{Feature Scaling}}: Menyeragamkan skala jangkauan nilai fitur numerik melalui standardisasi (\textit{Z-score scaling}) atau normalisasi (\textit{Min-Max scaling}) guna mencegah variabel bernilai besar mendominasi fungsi objektif model.
    \item \textbf{\textit{Dataset Splitting}}: Membagi himpunan data menjadi data latih (\textit{training set}) dan data uji (\textit{testing set}), misalnya dengan rasio 80:20, guna menguji daya generalisasi model secara objektif.
\end{enumerate}
Seluruh alur prapemrosesan ini diimplementasikan secara efisien menggunakan ekosistem pustaka komputasi Python, terutama Pandas, NumPy, dan Scikit-Learn [\ref{mckinney}].
\begin{figure}[H]
    \centering
    \includegraphics[width=0.6\textwidth]{62.png}
    \caption{Alur tahapan persiapan data dalam Data Preprocessing}
    \label{fig:62}
\end{figure}

\chapter{Hasil dan Pembahasan}

\section{Hasil Praktikum}

% ==========================================
% EXPLORATORY DATA ANALYSIS (EDA)
% ==========================================

\subsection{Exploratory Data Analysis (EDA)}

\subsubsection{1. Impor Pustaka}
\begin{figure}[H]
    \centering
    \includegraphics[width=0.6\textwidth]{1.png}
    \caption{Source Code Impor Pustaka}
    \label{fig:1}
\end{figure}
Penjelasan Gambar :
\begin{itemize}
    \item \textbf{Input:} Mengimpor pustaka yang dibutuhkan untuk manipulasi data dan visualisasi visual, yaitu \texttt{pandas} (\texttt{pd}), \texttt{numpy} (\texttt{np}), \texttt{matplotlib.pyplot} (\texttt{plt}), dan \texttt{seaborn} (\texttt{sns}).
\end{itemize}

\subsubsection{2. Load dataset}
\begin{figure}[H]
    \centering
    \includegraphics[width=0.6\textwidth]{2.png}
    \caption{Source Code Load dataset}
    \label{fig:2}
\end{figure}

\begin{figure}[H]
    \centering
    \includegraphics[width=0.8\textwidth]{3.png}
    \caption{Output Load dataset}
    \label{fig:3}
\end{figure}
Penjelasan Gambar :
\begin{itemize}
    \item \textbf{Input:} Menggunakan library \texttt{gdown} untuk mengunduh dua file CSV (\texttt{churning.csv} dan \texttt{new\_data.csv}) langsung dari Google Drive menggunakan file ID masing-masing. File tersebut kemudian dibaca menjadi objek DataFrame pandas (\texttt{df\_churning} dan \texttt{df\_new\_data}), dilanjutkan dengan perintah pencetakan informasi bentuk (shape) tabel.
    \item \textbf{Output:} Menampilkan log terminal proses pengunduhan hingga mencapai 100\%, diikuti oleh output pencetakan dimensi dataset. \texttt{df\_churning} memiliki ukuran (10127, 23) yang berarti 10.127 baris dan 23 kolom, sementara \texttt{df\_new\_data} berukuran (8143, 7).
\end{itemize}

\subsubsection{3. Melihat lima baris pertama dengan df.head()}
\begin{figure}[H]
    \centering
    \includegraphics[width=0.8\textwidth]{4.png}
    \caption{Source Code dan Output df.head()}
    \label{fig:4}
\end{figure}
Penjelasan Gambar :
\begin{itemize}
    \item \textbf{Input:} Mengeksekusi fungsi \texttt{df\_churning.head()} untuk melihat cuplikan sebagian kecil data teratas.
    \item \textbf{Output:} Menampilkan struktur tabel interaktif berisi 5 baris pertama data (indeks 0 hingga 4) lengkap dengan seluruh kolom variabelnya (seperti \texttt{CLIENTNUM}, \texttt{Attrition\_Flag}, dll).
\end{itemize}

\subsubsection{4. Melihat lima baris terakhir dengan df.tail()}
\begin{figure}[H]
    \centering
    \includegraphics[width=0.8\textwidth]{5.png}
    \caption{Source Code dan Output df.tail()}
    \label{fig:5}
\end{figure}
Penjelasan Gambar :
\begin{itemize}
    \item \textbf{Input:} Memanggil fungsi \texttt{df\_churning.tail()} untuk mengekstraksi dan meninjau cuplikan data paling bawah.
    \item \textbf{Output:} Menampilkan tabel yang memuat 5 baris data paling akhir dari DataFrame (mulai dari indeks 10122 hingga 10126).
\end{itemize}

\subsubsection{5. Menyimpan nilai dari class menjadi 2 dataframe yang berbeda}
\begin{figure}[H]
    \centering
    \includegraphics[width=0.8\textwidth]{6.png}
    \caption{Source Code Menyimpan nilai dari class ke DataFrame berbeda}
    \label{fig:6}
\end{figure}
Penjelasan Gambar :
\begin{itemize}
    \item \textbf{Input:} Menerapkan seleksi kondisi (\textit{boolean indexing}) pada \texttt{df\_churning}. Baris data dengan \texttt{Attrition\_Flag} bernilai 'Existing Customer' dipisah dan disimpan ke variabel \texttt{existing\_data}, sedangkan 'Attrited Customer' disimpan ke \texttt{churn\_data}.
\end{itemize}

\subsubsection{6. Line chart}
\begin{figure}[H]
    \centering
    \includegraphics[width=0.55\textwidth]{7.png}
    \caption{Source Code Line chart}
    \label{fig:7}
\end{figure}

\begin{figure}[H]
    \centering
    \includegraphics[width=0.55\textwidth]{8.png}
    \caption{Output Line chart}
    \label{fig:8}
\end{figure}
Penjelasan Gambar :
\begin{itemize}
    \item \textbf{Input:} 
    \begin{itemize}
        \item Blok pertama mendefinisikan array statis untuk membuat dua garis berbeda (\texttt{y1} bermarker 'o' dan \texttt{y2} bermarker 'x') di atas sumbu \texttt{x} yang sama menggunakan \texttt{plt.plot()}.
        \item Blok kedua memplot grafik deret waktu (time series) dari dataset IoT setelah mengonversi tipe kolom 'date' menjadi format \texttt{datetime}, dengan lebar figur disesuaikan menjadi 15x3.
    \end{itemize}
    \item \textbf{Output:}
    \begin{itemize}
        \item \textbf{Grafik 1 (Atas):} Menghasilkan \textit{line chart} dasar dengan garis biru (data y1) dan garis oranye (data y2) yang saling berpotongan terhadap titik x.
        \item \textbf{Grafik 2 (Bawah):} Menghasilkan \textit{line chart} memanjang yang memvisualisasikan tren fluktuasi kelembapan udara (\texttt{Humidity}) berdasarkan data waktu (\texttt{date}).
    \end{itemize}
\end{itemize}

\subsubsection{7. Pie chart}
\begin{figure}[H]
    \centering
    \includegraphics[width=0.50\textwidth]{9.png}
    \caption{Source Code Pie chart}
    \label{fig:9}
\end{figure}

\begin{figure}[H]
    \centering
    \includegraphics[width=0.30\textwidth]{10.png}
    \caption{Output Pie chart}
    \label{fig:10}
\end{figure}
Penjelasan Gambar :
\begin{itemize}
    \item \textbf{Input:}
    \begin{itemize}
        \item Bagian atas membangun \textit{pie chart} menggunakan \texttt{plt.pie()} dari sekumpulan nilai kuantitas \texttt{y} yang diinput manual beserta daftar label spesifiknya.
        \item Bagian bawah secara dinamis menciptakan \textit{pie chart} dengan menghitung agregat kemunculan tipe status pernikahan dari kolom \texttt{Marital\_Status} menggunakan fungsi \texttt{value\_counts()}, lalu mengekstrak indeksnya untuk dijadikan label diagram.
    \end{itemize}
    \item \textbf{Output:}
    \begin{itemize}
        \item \textbf{Grafik 1 (Atas):} Menampilkan diagram lingkaran distribusi profesi yang terbagi menjadi Web Programmer (35.00\%), Data Scientist (25.00\%), DB Admin (25.00\%), dan Manager (15.00\%).
        \item \textbf{Grafik 2 (Bawah):} Menampilkan diagram proporsi \textit{Marital Status} keseluruhan pelanggan, didominasi oleh kategori Married (46.28\%) dan Single (38.94\%).
    \end{itemize}
\end{itemize}

\subsubsection{8. Bar plot}
\begin{figure}[H]
    \centering
    \includegraphics[width=0.6\textwidth]{11.png}
    \caption{Source Code Bar plot}
    \label{fig:11}
\end{figure}

\begin{figure}[H]
    \centering
    \includegraphics[width=0.6\textwidth]{12.png}
    \caption{Output Bar plot}
    \label{fig:12}
\end{figure}
Penjelasan Gambar :
\begin{itemize}
    \item \textbf{Input:} 
    \begin{itemize}
        \item Bagian atas membuat diagram batang (bar plot) sederhana menggunakan \texttt{plt.bar(x,y)} dengan data \textit{dummy} profesi dan memutar label sumbu X sebesar 45 derajat.
        \item Bagian bawah membuat diagram batang berdampingan (\textit{grouped bar chart}) dengan menghitung jumlah pelanggan berdasarkan \texttt{Marital\_Status} untuk kelompok \textit{Existing} dan \textit{Churned}, menyamakan indeks kategorinya, dan memplotnya menggunakan pengaturan \texttt{bar\_width} (lebar batang) dan penyesuaian posisi indeks.
    \end{itemize}
    \item \textbf{Output:} 
    \begin{itemize}
        \item \textbf{Grafik 1 (Atas):} Menampilkan diagram batang sederhana yang memvisualisasikan jumlah untuk setiap profesi (Web Programmer, Data Scientist, DB Admin, Manager).
        \item \textbf{Grafik 2 (Bawah):} Menampilkan \textit{grouped bar chart} yang membandingkan secara langsung jumlah pelanggan aktif (biru) dan pelanggan yang berhenti/churn (oranye) pada tiap kategori status pernikahan.
    \end{itemize}
\end{itemize}

\subsubsection{9. Stacked bar plot}
\begin{figure}[H]
    \centering
    \includegraphics[width=0.6\textwidth]{13.png}
    \caption{Source Code Stacked bar plot}
    \label{fig:13}
\end{figure}

\begin{figure}[H]
    \centering
    \includegraphics[width=0.6\textwidth]{14.png}
    \caption{Output Stacked bar plot}
    \label{fig:14}
\end{figure}
Penjelasan Gambar :
\begin{itemize}
    \item \textbf{Input:} Menghitung kemunculan tiap kategori \texttt{Marital\_Status} pada data \textit{existing} dan \textit{churn}, lalu menggabungkannya ke dalam satu DataFrame baru menggunakan \texttt{pd.concat()} berdasarkan kolom (\texttt{axis=1}). Plot dibuat dengan memanggil metode \texttt{.plot(kind='bar', stacked=True)}.
    \item \textbf{Output:} Menampilkan diagram batang bertumpuk (\textit{stacked bar chart}) di mana jumlah pelanggan \textit{Churned} (oranye) ditumpuk di atas pelanggan \textit{Existing} (biru) untuk setiap kategori status pernikahan, memberikan gambaran proporsi secara visual.
\end{itemize}

\subsubsection{10. Histogram}
\begin{figure}[H]
    \centering
    \includegraphics[width=0.55\textwidth]{15.png}
    \caption{Source Code Histogram}
    \label{fig:14}
\end{figure}

\begin{figure}[H]
    \centering
    \includegraphics[width=0.35\textwidth]{16.png}
    \caption{Output Histogram}
    \label{fig:16}
\end{figure}
Penjelasan Gambar :
\begin{itemize}
    \item \textbf{Input:} 
    \begin{itemize}
        \item Bagian atas membuat histogram dari kumpulan angka acak dengan mendefinisikan batas kelas interval (bins) secara manual yaitu \texttt{[0, 25, 50, 75, 100]}.
        \item Bagian bawah membuat histogram yang menumpuk/menggabungkan (overlay) distribusi \texttt{Customer\_Age} antara pelanggan \textit{Existing} dan pelanggan \textit{Churned} pada satu sumbu (Axes) yang sama.
    \end{itemize}
    \item \textbf{Output:} 
    \begin{itemize}
        \item \textbf{Grafik 1 (Atas):} Menampilkan frekuensi nilai berdasarkan 4 interval kelas (0-25, 25-50, 50-75, 75-100).
        \item \textbf{Grafik 2 (Bawah):} Menampilkan persebaran usia pelanggan (\textit{Customer Age}), di mana distribusi umur untuk pelanggan yang bertahan (biru) menumpuk bersama distribusi pelanggan yang berhenti (oranye).
    \end{itemize}
\end{itemize}

\subsubsection{11. Scatter plot}
\begin{figure}[H]
    \centering
    \includegraphics[width=0.45\textwidth]{17.png}
    \caption{Source Code Scatter plot}
    \label{fig:17}
\end{figure}

\begin{figure}[H]
    \centering
    \includegraphics[width=0.35\textwidth]{18.png}
    \caption{Output Scatter plot}
    \label{fig:18}
\end{figure}
Penjelasan Gambar :
\begin{itemize}
    \item \textbf{Input:} 
    \begin{itemize}
        \item \textbf{Versi 1 (Atas):} Menggunakan fungsi \texttt{plt.scatter()} dari library Matplotlib untuk memplot titik koordinat Usia versus Lama Menjadi Nasabah (\textit{Months on book}), dipanggil dua kali berturut-turut untuk grup yang berbeda.
        \item \textbf{Versi 2 (Bawah):} Menggunakan metode bawaan Pandas yakni \texttt{.plot.scatter()}, di mana objek grafik (\textit{Axes}) dari plot pertama dilempar ke argumen \texttt{ax=ax} pada plot kedua agar titik dirender pada kanvas yang sama, dengan warna yang disetel manual (merah dan biru).
    \end{itemize}
    \item \textbf{Output:} 
    \begin{itemize}
        \item \textbf{Grafik 1 (Atas):} \textit{Scatter plot} standar dengan warna default matplotlib (biru dan oranye) yang menunjukkan korelasi usia terhadap lama berlangganan.
        \item \textbf{Grafik 2 (Bawah):} \textit{Scatter plot} yang menghasilkan relasi bentuk yang sama, namun menggunakan kustomisasi warna khusus, yaitu biru (\textit{churn}) dan merah (\textit{existing}).
    \end{itemize}
\end{itemize}

\subsubsection{12. Box plot}
\begin{figure}[H]
    \centering
    \includegraphics[width=0.45\textwidth]{19.png}
    \caption{Source Code Box plot}
    \label{fig:19}
\end{figure}

\begin{figure}[H]
    \centering
    \includegraphics[width=0.35\textwidth]{20.png}
    \caption{Output Box plot}
    \label{fig:20}
\end{figure}
Penjelasan Gambar :
\begin{itemize}
    \item \textbf{Input:} 
    \begin{itemize}
        \item \textbf{Versi 1 (Atas):} Menggunakan \texttt{plt.boxplot()} untuk menghasilkan analisis bentuk distribusi dasar dari satu kolom penuh yakni \texttt{Customer\_Age} pada seluruh dataset gabungan.
        \item \textbf{Versi 2 (Bawah):} Menggabungkan kolom usia dari pelanggan \textit{Existing} dan \textit{Churned} menjadi struktur data baru yang sejajar, lalu memanggil fungsi \texttt{df\_new.plot(kind='box')} melalui Pandas.
    \end{itemize}
    \item \textbf{Output:} 
    \begin{itemize}
        \item \textbf{Grafik 1 (Atas):} \textit{Box plot} tunggal yang menampilkan nilai median, kuartil bawah/atas, nilai ekstrim (whisker), dan mendeteksi adanya data pencilan (outlier) berupa titik lingkaran di bagian atas.
        \item \textbf{Grafik 2 (Bawah):} Menghasilkan dua \textit{box plot} berdampingan yang mempermudah perbandingan statistik persebaran umur antara kelompok nasabah lama (\textit{Existing}) dengan nasabah yang keluar (\textit{Churned}).
    \end{itemize}
\end{itemize}

\subsubsection{13. Violin plot}
\begin{figure}[H]
    \centering
    \includegraphics[width=0.6\textwidth]{21.png}
    \caption{Source Code Violin plot}
    \label{fig:21}
\end{figure}

\begin{figure}[H]
    \centering
    \includegraphics[width=0.6\textwidth]{22.png}
    \caption{Output Violin plot}
    \label{fig:22}
\end{figure}
Penjelasan Gambar :
\begin{itemize}
    \item \textbf{Input:} Mengekstrak nilai usia (\textit{Customer\_Age}) dari nasabah \textit{Existing} dan \textit{Churned} ke dalam bentuk array, kemudian menyimpannya ke dalam struktur \textit{dictionary}. Fungsi \texttt{plt.violinplot()} dipanggil untuk memplot nilai-nilai tersebut, dan \texttt{plt.xticks()} digunakan untuk mengganti label sumbu X sesuai dengan kunci (keys) dari dictionary.
    \item \textbf{Output:} Menghasilkan diagram biola (\textit{violin plot}) yang menunjukkan distribusi kepadatan probabilitas dari usia nasabah. Bentuk yang menggelembung menunjukkan konsentrasi data terbanyak, di mana kedua kategori memiliki kepadatan tertinggi pada rentang usia 40 hingga 50 tahun.
\end{itemize}

\subsubsection{14. Sub plot}
\begin{figure}[H]
    \centering
    \includegraphics[width=0.45\textwidth]{23.png}
    \caption{Source Code Sub plot}
    \label{fig:23}
\end{figure}

\begin{figure}[H]
    \centering
    \includegraphics[width=0.45\textwidth]{24.png}
    \caption{Output Sub plot}
    \label{fig:24}
\end{figure}
Penjelasan Gambar :
\begin{itemize}
    \item \textbf{Input:} Membagi satu area kanvas figure (ukuran 10x10) menjadi grid 2x2 menggunakan \texttt{plt.subplot(baris, kolom, posisi)}. Posisi 1 dan 2 diisi dengan diagram kotak (\textit{boxplot}) untuk usia masing-masing kategori, sedangkan posisi 3 dan 4 diisi dengan histogram usia, dengan penyelarasan batas sumbu Y menggunakan \texttt{plt.ylim([0, 2100])}.
    \item \textbf{Output:} Menampilkan empat grafik independen dalam satu kesatuan gambar (kanvas). Baris atas menampilkan perbandingan \textit{boxplot} usia, dan baris bawah menampilkan \textit{histogram} jumlah rentang usia dari pelanggan lama (biru) dan pelanggan yang berhenti (merah).
\end{itemize}

\subsubsection{15. Annotation}
\begin{figure}[H]
    \centering
    \includegraphics[width=0.6\textwidth]{25.png}
    \caption{Source Code Annotation}
    \label{fig:25}
\end{figure}

\begin{figure}[H]
    \centering
    \includegraphics[width=0.65\textwidth]{26.png}
    \caption{Output Annotation}
    \label{fig:26}
\end{figure}
Penjelasan Gambar :
\begin{itemize}
    \item \textbf{Input:} Membuat histogram dari seluruh data usia pelanggan. Setelah itu, fungsi \texttt{plt.text(x, y, "teks")} digunakan untuk menyematkan anotasi berupa teks tambahan ke dalam koordinat grafik secara manual, yang berisi kalkulasi dinamis nilai Median dan Mean dari data tersebut.
    \item \textbf{Output:} Menghasilkan visualisasi histogram yang dilengkapi dengan teks informasi stasistik tambahan (Mean Age: 46.33 dan Median Age: 46.0) yang melayang di koordinat ruang kosong sebelah kanan atas grafik.
\end{itemize}

\subsubsection{16. Axis}
\begin{figure}[H]
    \centering
    \includegraphics[width=0.6\textwidth]{27.png}
    \caption{Source Code Axis}
    \label{fig:27}
\end{figure}

\begin{figure}[H]
    \centering
    \includegraphics[width=0.6\textwidth]{28.png}
    \caption{Output Axis}
    \label{fig:28}
\end{figure}
Penjelasan Gambar :
\begin{itemize}
    \item \textbf{Input:} Membuat grafik garis (\textit{line chart}) ganda sederhana. Pengaturan \textit{axis} (sumbu) dikonfigurasi melalui fungsi \texttt{plt.ylim([-20, 30])} untuk memaksakan batas minimum dan maksimum sumbu Y, serta mengubah ukuran \textit{font} label pada sumbu X dan Y menjadi lebih besar (\texttt{fontsize=18}).
    \item \textbf{Output:} Menampilkan grafik garis dengan ruang kosong yang cukup besar di bagian bawah, karena rentang nilai sumbu Y telah dimodifikasi (di-skala) secara manual agar dimulai dari angka negatif (-20) hingga angka batas 30.
\end{itemize}

\subsubsection{17. Legend}
\begin{figure}[H]
    \centering
    \includegraphics[width=0.50\textwidth]{29.png}
    \caption{Source Code Legend}
    \label{fig:29}
\end{figure}

\begin{figure}[H]
    \centering
    \includegraphics[width=0.50\textwidth]{30.png}
    \caption{Output Legend}
    \label{fig:30}
\end{figure}
Penjelasan Gambar :
\begin{itemize}
    \item \textbf{Input:} Memplot dua histogram (distribusi data usia \textit{Existing} dan \textit{Churned}) dalam satu sumbu yang sama. Argumen \texttt{label='nama\_kategori'} diberikan saat pemanggilan \texttt{plt.hist()}, kemudian dimunculkan secara visual melalui pemanggilan metode \texttt{plt.legend()}.
    \item \textbf{Output:} Menampilkan histogram bertumpuk yang dilengkapi dengan kotak legenda di sudut kanan atas. Kotak ini memberikan indikator informasi warna secara visual (biru untuk Existing Customer dan oranye untuk Churned Customer), sehingga mempermudah pembacaan grafik multi-data.
\end{itemize}


% ==========================================
% DATA PREPROCESSING
% ==========================================
\clearpage %
\subsection{Data Preprocessing}

\subsubsection{1. Impor pustaka}
\begin{figure}[H]
    \centering
    \includegraphics[width=0.6\textwidth]{31.png}
    \caption{Source Code Impor pustaka}
    \label{fig:31}
\end{figure}
Penjelasan Gambar :
\begin{itemize}
    \item \textbf{Input:} Memuat semua library esensial untuk manipulasi data (\texttt{pandas}), plotting visual (\texttt{matplotlib}), pengunduhan file eksternal (\texttt{gdown}), serta berbagai modul spesifik dari \texttt{scikit-learn} (\texttt{sklearn}) yang diperlukan dalam tahap \textit{preprocessing} (\texttt{train\_test\_split}, \texttt{OneHotEncoder}, \texttt{LabelEncoder}, dll).
\end{itemize}

\subsubsection{2. Load dataset}
\begin{figure}[H]
    \centering
    \includegraphics[width=0.6\textwidth]{32.png}
    \caption{Source Code Load dataset}
    \label{fig:32}
\end{figure}
Penjelasan Gambar :
\begin{itemize}
    \item \textbf{Input:} Mengeksekusi pengunduhan file dataset \textit{Titanic} dari Google Drive menggunakan \texttt{gdown} melalui ID unik tautan berbaginya, kemudian menyimpannya dengan nama \texttt{data\_from\_drive.csv}, lalu langsung membaca file tersebut ke dalam objek DataFrame bernama \texttt{df}. Proses pengunduhan juga memunculkan pesan status keberhasilan (100\%).
\end{itemize}

\subsubsection{3. Melihat lima baris data pertama dengan df.head()}
\begin{figure}[H]
    \centering
    \includegraphics[width=0.8\textwidth]{33.png}
    \caption{Source Code dan Output df.head()}
    \label{fig:33}
\end{figure}
Penjelasan Gambar :
\begin{itemize}
    \item \textbf{Input:} Menjalankan perintah \texttt{df.head()} untuk mengintip format tabel bagian atas.
    \item \textbf{Output:} Mencetak lima baris paling awal (indeks 0 hingga 4) dari dataset kapal Titanic, memperlihatkan rincian kolom penting seperti \texttt{Survived} (keselamatan), \texttt{Pclass} (kelas tiket), \texttt{Name}, \texttt{Sex}, dan atribut lainnya.
\end{itemize}

\subsubsection{4. Melihat lima baris data terkahir dengan df.tail()}
\begin{figure}[H]
    \centering
    \includegraphics[width=0.8\textwidth]{34.png}
    \caption{Source Code dan Output df.tail()}
    \label{fig:34}
\end{figure}
Penjelasan Gambar :
\begin{itemize}
    \item \textbf{Input:} Menjalankan perintah \texttt{df.tail()} untuk mengekstraksi dan meninjau cuplikan data dari ujung paling bawah tabel.
    \item \textbf{Output:} Menampilkan struktur baris penutup dari DataFrame (indeks 886 hingga 890), yang secara implisit juga memberikan petunjuk tentang total keseluruhan observasi dalam dataset tersebut.
\end{itemize}

\subsubsection{5. Melihat bentuk dataset dengan df.shape}
\begin{figure}[H]
    \centering
    \includegraphics[width=0.6\textwidth]{35.png}
    \caption{Source Code dan Output Melihat bentuk dataset dengan df.shape}
    \label{fig:35}
\end{figure}
Penjelasan Gambar :
\begin{itemize}
    \item \textbf{Input:} Memanggil properti \texttt{.shape} pada DataFrame \texttt{df} untuk mengekstrak informasi struktural dimensi data.
    \item \textbf{Output:} Mengembalikan nilai \textit{tuple} \texttt{(891, 12)}, yang menegaskan bahwa tabel dataset Titanic yang diunduh terdiri dari tepat 891 baris (penumpang) dan 12 kolom (atribut).
\end{itemize}

\subsubsection{6. Visualisasi distribusi data}
\begin{figure}[H]
    \centering
    \includegraphics[width=0.6\textwidth]{36.png}
    \caption{Source Code Visualisasi distribusi data}
    \label{fig:36}
\end{figure}

\begin{figure}[H]
    \centering
    \includegraphics[width=0.6\textwidth]{37.png}
    \caption{Output Visualisasi distribusi data}
    \label{fig:37}
\end{figure}
Penjelasan Gambar :
\begin{itemize}
    \item \textbf{Input:} 
    \begin{itemize}
        \item Blok pertama membuat \textit{pie chart} biner (dua potongan) untuk memetakan distribusi kolom target klasifikasi, yaitu proporsi jumlah penumpang yang tewas (0) vs selamat (1) dari kolom \texttt{Survived}.
        \item Blok kedua membuat kumpulan \textit{histogram} serentak untuk seluruh variabel dengan tipe data numerik yang ada di dalam \textit{dataframe}, menggunakan metode iterasi bawaan \texttt{df.hist()} dengan luas kanvas 10x10.
    \end{itemize}
    \item \textbf{Output:} 
    \begin{itemize}
        \item \textbf{Grafik 1 (Atas):} Menghasilkan diagram lingkaran porsi keselamatan, di mana persentase jumlah korban meninggal (label 0, warna biru) secara tragis mendominasi jauh di atas persentase penumpang yang berhasil selamat (label 1, warna oranye).
        \item \textbf{Grafik 2 (Bawah):} Menghasilkan matriks grid berisi deretan \textit{histogram} kecil. Ini memudahkan analis untuk menginspeksi kurva distribusi nilai variabel numerikal krusial seperti umur (\texttt{Age}), harga tiket (\texttt{Fare}), status keluarga (\texttt{SibSp}, \texttt{Parch}), dll secara bersamaan.
    \end{itemize}
\end{itemize}

\subsubsection{7. Analisis korelasi antar variabel}
\begin{figure}[H]
    \centering
    \includegraphics[width=0.8\textwidth]{38.png}
    \caption{Source Code dan Output Analisis korelasi antar variabel}
    \label{fig:38}
\end{figure}
Penjelasan Gambar :
\begin{itemize}
    \item \textbf{Input:} Mengeksekusi fungsi analitik statistik \texttt{df.corr()} dengan parameter pelindung \texttt{numeric\_only=True} (untuk mengabaikan kolom berformat \textit{string/teks}) guna mengkalkulasi koefisien korelasi linier Pearson antar semua atribut kuantitatif yang ada.
    \item \textbf{Output:} Mencetak matriks tabel dua dimensi yang menjabarkan tingkat korelasi silang. Nilai 1 merepresentasikan korelasi absolut pada variabel yang sama (diagonal). Nilai mendekati 1 atau -1 mengindikasikan adanya hubungan yang kuat, sementara nilai mendekati 0 menandakan lemahnya korelasi antar dua fitur (misal hubungan kuat negatif antara \texttt{Pclass} dan \texttt{Fare} sebesar -0.54).
\end{itemize}

\subsubsection{8. Descriptive statistics}
\begin{figure}[H]
    \centering
    \includegraphics[width=0.8\textwidth]{39.png}
    \caption{Source Code dan Output Descriptive statistics}
    \label{fig:39}
\end{figure}
Penjelasan Gambar :
\begin{itemize}
    \item \textbf{Input:} Memanggil perintah \texttt{df.describe()} untuk menghasilkan rekapitulasi hitungan statistik univariat dasar terhadap seluruh variabel bertipe numerik.
    \item \textbf{Output:} Menghasilkan tabel deskriptif padat yang membedah ukuran pemusatan dan penyebaran tiap kolom. Tabel ini menginformasikan metrik jumlah data valid (\textit{count}), nilai rata-rata (\textit{mean}), standar deviasi (\textit{std}), batas nilai ekstrim bawah (\textit{min}) dan atas (\textit{max}), serta pemisahan distribusi data melalui nilai persentil kuartil (25\%, 50\%/median, dan 75\%).
\end{itemize}

\subsubsection{9. Mengecek missing values dan duplicates}
\begin{figure}[H]
    \centering
    \includegraphics[width=0.6\textwidth]{40.png}
    \caption{Source Code Mengecek missing values dan duplicates}
    \label{fig:40}
\end{figure}
Penjelasan Gambar :
\begin{itemize}
    \item \textbf{Input:} 
    \begin{itemize}
        \item Bagian pertama mengecek keberadaan data sel yang kosong (\textit{null/NaN}) di tiap kolom menggunakan fungsi kombinasi \texttt{df.isnull().sum()}.
        \item Bagian kedua melakukan deteksi perulangan baris data identik menggunakan \texttt{df.duplicated().sum()}, dan memusnahkan duplikasi tesebut secara permanen pada memori data menggunakan \texttt{df.drop\_duplicates(inplace=True)}.
    \end{itemize}
\end{itemize}

\subsubsection{10. Imputation}
\begin{figure}[H]
    \centering
    \includegraphics[width=0.6\textwidth]{41.png}
    \caption{Source Code Imputation}
    \label{fig:41}
\end{figure}

\begin{figure}[H]
    \centering
    \includegraphics[width=0.8\textwidth]{42.png}
    \caption{Output Imputation}
    \label{fig:42}
\end{figure}
Penjelasan Gambar :
\begin{itemize}
    \item \textbf{Input:} Melakukan penanganan nilai kosong (\textit{missing value imputation}). Untuk kolom numerik umur (\texttt{Age}), kekosongan diisi dengan nilai median menggunakan \texttt{fillna()}. Sementara untuk kolom bertipe data kategorikal (\texttt{Cabin} dan \texttt{Embarked}), sel yang kosong ditambal dengan nilai modus atau nilai yang paling sering muncul yang didapatkan dari pemanggilan \texttt{value\_counts().index[0]}.
    \item \textbf{Output:} Menghasilkan pesan \textit{FutureWarning} karena penggunaan parameter modifikasi \texttt{inplace=True} pada Pandas versi terbaru. Namun yang terpenting, cetakan dari \texttt{df.isnull().sum()} di bagian bawah menegaskan seluruh kolom telah menampilkan angka 0. Ini menandakan bahwa proses \textit{imputation} berhasil membersihkan tabel dari semua elemen data kosong.
\end{itemize}

\subsubsection{11. Memeriksa categorical attribute}
\begin{figure}[H]
    \centering
    \includegraphics[width=0.8\textwidth]{43.png}
    \caption{Source Code dan Output Memeriksa categorical attribute}
    \label{fig:43}
\end{figure}
Penjelasan Gambar :
\begin{itemize}
    \item \textbf{Input:} Memisahkan dataframe menjadi himpunan fitur independen (\texttt{df\_x}) dan dependen (\texttt{df\_y}). Kemudian, mendeteksi dan mengekstrak nama kolom yang memiliki tipe data objek/string kategorikal menggunakan \texttt{select\_dtypes(include=['object', 'bool']).columns}.
    \item \textbf{Output:} Mencetak daftar indeks nama kolom kategorikal yang ditemukan di dalam dataset, yaitu kolom \texttt{Sex} dan \texttt{Embarked}.
\end{itemize}

\subsubsection{12. One hot encoding}
\begin{figure}[H]
    \centering
    \includegraphics[width=0.6\textwidth]{44.png}
    \caption{Source Code dan Output One hot encoding}
    \label{fig:44}
\end{figure}
Penjelasan Gambar :
\begin{itemize}
    \item \textbf{Input:} Menerapkan metode transformasi \textit{One-Hot Encoding} pada kolom kategorikal untuk mengkonversi nilai teks menjadi matriks fitur biner. Gambar menyajikan dua teknik berbeda: pertama, menggunakan fungsi bawaan Pandas yakni \texttt{pd.get\_dummies()}; kedua, menggunakan fungsi \textit{machine learning} \texttt{OneHotEncoder()} dari pustaka \textit{sklearn}.
    \item \textbf{Output:} 
    \begin{itemize}
        \item Cara Pandas menghasilkan tabel \textit{dataframe} baru dengan kolom perluasan (seperti \texttt{Sex\_male}, \texttt{Embarked\_Q}, \texttt{Embarked\_S}), di mana keberadaan sifat ditandai dengan nilai boolean \textit{True/False}.
        \item Cara Scikit-Learn mereturn representasi data berupa matriks \textit{array} numpy mentah yang berisi susunan sekuensial deret angka desimal \texttt{1.0} dan \texttt{0.0}.
    \end{itemize}
\end{itemize}

\subsubsection{13. Label encoding}
\begin{figure}[H]
    \centering
    \includegraphics[width=0.6\textwidth]{45.png}
    \caption{Source Code dan Output Label encoding fitur}
    \label{fig:45}
\end{figure}

\begin{figure}[H]
    \centering
    \includegraphics[width=0.6\textwidth]{46.png}
    \caption{Source Code dan Output Label encoding target}
    \label{fig:46}
\end{figure}
Penjelasan Gambar :
\begin{itemize}
    \item \textbf{Input:} Memanfaatkan teknik alternatif encoding untuk mengubah kategori teks menjadi nilai nominal terurut numerik via \textit{Label Encoding}. 
    \begin{itemize}
        \item Gambar pertama menggunakan fungsi iterasi (\textit{looping for}) untuk mengimplementasikan \texttt{LabelEncoder().fit\_transform()} pada semua kolom fitur independen (\texttt{df\_x}) yang bertipe kategorikal.
        \item Gambar kedua secara spesifik menerapkan \texttt{LabelEncoder()} pada kolom tunggal label target prediksi keselamatan (\texttt{df\_y}).
    \end{itemize}
    \item \textbf{Output:} 
    \begin{itemize}
        \item \textbf{Grafik 1 (Atas):} Nilai teks pada kolom \texttt{Sex} terganti dengan kelas angka (0 untuk \textit{female}, 1 untuk \textit{male}), serta nilai pada kolom \texttt{Embarked} dikonversi menjadi angka pemetaan (0, 1, 2). 
        \item \textbf{Grafik 2 (Bawah):} Transformasi label merubah nilai pada kolom target prediksi menjadi sebuah susunan \textit{array} flat satu dimensi yang murni berisi angka 1 (selamat) dan 0 (tidak selamat).
    \end{itemize}
\end{itemize}

\subsubsection{14. Mengecek korelasi atribut}
\begin{figure}[H]
    \centering
    \includegraphics[width=0.8\textwidth]{47.png}
    \caption{Source Code dan Output Mengecek korelasi atribut}
    \label{fig:47}
\end{figure}
Penjelasan Gambar :
\begin{itemize}
    \item \textbf{Input:} Menyatukan kembali variabel target \texttt{df\_y} ke dalam himpunan fitur prediktor \texttt{df\_x} menggunakan fungsi \texttt{pd.concat(axis=1)}. Penggabungan ini diperlukan untuk menjalankan kalkulasi silang fungsi korelasi Pearson komprehensif via \texttt{df\_gabung.corr()}.
    \item \textbf{Output:} Mengembalikan \textit{dataframe} matriks koefisien. Karena kolom-kolom string seperti \texttt{Sex} dan \texttt{Embarked} telah direpresentasikan menjadi angka pada tahap sebelumnya, kini tingkat hubungan linear korelasinya terhadap probabilitas \texttt{Survived} dapat dihitung dan ditampilkan secara utuh.
\end{itemize}

\subsubsection{15. Membagi dataset}
\begin{figure}[H]
    \centering
    \includegraphics[width=0.8\textwidth]{48.png}
    \caption{Source Code dan Output Membagi dataset}
    \label{fig:48}
\end{figure}
Penjelasan Gambar :
\begin{itemize}
    \item \textbf{Input:} Menjalankan skema \textit{Data Splitting} menggunakan fungsi \texttt{train\_test\_split()}. Proporsi ukuran pembagian diatur pada argumen \texttt{test\_size=0.3} (berarti 30\% data dicadangkan sebagai \textit{Test Set} dan 70\% menjadi \textit{Training Set}). Argumen \texttt{random\_state=42} dipasang agar pembagian acak ini bersifat konsisten saat dijalankan ulang.
    \item \textbf{Output:} Menampilkan potongan tabel fitur latih (\texttt{X\_train}). Akibat pemotongan proporsional 70\%, tabel \texttt{X\_train} kini memiliki bentuk dimensi yang menyusut menjadi 623 baris data, dengan posisi urutan baris indeks yang telah teracak.
\end{itemize}

\subsubsection{16. Scaling}
\begin{figure}[H]
    \centering
    \includegraphics[width=0.6\textwidth]{49.png}
    \caption{Source Code dan Output Scaling}
    \label{fig:49}
\end{figure}
Penjelasan Gambar :
\begin{itemize}
    \item \textbf{Input:} Mengaplikasikan \textit{Feature Scaling} pada dataset latih dan uji menggunakan dua pendekatan transformer dari \textit{sklearn} secara bergantian.
    \begin{itemize}
        \item Blok atas menerapkan \texttt{StandardScaler()} untuk memusatkan distribusi data (mean=0, varians=1).
        \item Blok bawah menerapkan \texttt{MinMaxScaler()} untuk memampatkan nilai ke dalam rentang rasio [0, 1].
    \end{itemize}
    \item \textbf{Output:} Kedua transformer menghasilkan susunan matriks Numpy termodifikasi. \textit{StandardScaler} memuat penyebaran angka desimal yang mencakup rentang nilai negatif dan positif, sedangkan \textit{MinMaxScaler} berhasil menekan seluruh elemen observasi absolut berada di antara batas angka 0 hingga batas angka 1.
\end{itemize}

% ==========================================
% TUGAS PRAKTIKUM
% ==========================================
\clearpage
\section{Tugas Praktikum}

\subsection{Exploratory Data Analysis (EDA)}

\begin{itemize}
    \item Dataset: https://www.kaggle.com/datasets/sakshigoyal7/credit-card-customers
    \item Silahkan membaca dataset diatas menggunakan dataframe!
    \item Jelaskan apa tujuan dari penggunaan dataset ini!
    \item Definisikan atribut mana yang menjadi input dan atribut mana yang menjadi output/class/label!
    \item Berikan penjelasan singkat untuk setiap atribut/variable dari dataset tersebut
\end{itemize}

\subsubsection{1. Impor Pustaka, Load Dataset, dan Pembersihan Awal}
\begin{figure}[H]
    \centering
    \includegraphics[width=0.8\textwidth]{50.png}
    \caption{Source Code Impor Pustaka dan Load Dataset}
    \label{fig:50}
\end{figure}
Penjelasan Gambar :
\begin{itemize}
    \item \textbf{Input:} Memuat pustaka manipulasi data (\texttt{pandas}, \texttt{numpy}) dan visualisasi (\texttt{matplotlib}, \texttt{seaborn}). Mengunduh dataset melalui \texttt{gdown} dan membacanya menjadi \textit{dataframe}. Kolom yang tidak relevan (seperti ID dan \textit{classifier} bawaan) langsung dihapus menggunakan \texttt{drop()}.
\end{itemize}

\subsubsection{2. Menampilkan Data Teratas dan Terbawah}
\begin{figure}[H]
    \centering
    \includegraphics[width=0.8\textwidth]{51.png}
    \caption{Output df\_bersih.head()}
    \label{fig:51}
\end{figure}

\begin{figure}[H]
    \centering
    \includegraphics[width=0.8\textwidth]{52.png}
    \caption{Output df\_bersih.tail()}
    \label{fig:52}
\end{figure}
Penjelasan Gambar :
\begin{itemize}
    \item \textbf{Input:} Mengeksekusi fungsi \texttt{head()} dan \texttt{tail()} pada tabel \textit{dataframe} yang telah dibersihkan.
    \item \textbf{Output:} Mencetak 5 baris pertama (indeks 0-4) dan 5 baris terakhir dari dataset \textit{Credit Card Customers}, memverifikasi bahwa kolom \texttt{Attrition\_Flag} dan fitur numerik lainnya sudah tersusun rapi untuk dianalisis.
\end{itemize}

\subsubsection{3. Analisis Dataset}
\textbf{Tujuan Penggunaan Dataset}
\begin{itemize}
    \item Dataset \textit{Credit Card Customers} ini digunakan untuk menganalisis dan memprediksi probabilitas nasabah kartu kredit yang akan berhenti menggunakan layanan bank (\textit{churn}). 
    \item Melalui \textit{Exploratory Data Analysis} (EDA), bank dapat memahami faktor-faktor pendorong utama keputusan nasabah untuk keluar, sehingga dapat merancang strategi retensi yang proaktif untuk mempertahankan nasabah tersebut.
\end{itemize}

\textbf{Definisi Atribut (Input dan Output)}
\begin{itemize}
    \item \textbf{Output / Class / Label:} Atribut \texttt{Attrition\_Flag} bertindak sebagai target prediksi (label) yang membedakan apakah nasabah bertahan (\textit{Existing Customer}) atau keluar (\textit{Attrited Customer}).
    \item \textbf{Input / Features:} Seluruh atribut lainnya di dalam dataset (seperti usia, batas kredit, status pernikahan, dll.) berfungsi sebagai variabel independen (input), dengan mengecualikan kolom \texttt{CLIENTNUM} dan kolom bawaan \textit{Naive Bayes Classifier}.
\end{itemize}

\textbf{Penjelasan Singkat Atribut Utama}
\begin{itemize}
    \item \texttt{Customer\_Age}: Usia nasabah dalam hitungan tahun.
    \item \texttt{Gender}: Jenis kelamin nasabah (\texttt{M} untuk pria, \texttt{F} untuk wanita).
    \item \texttt{Dependent\_count}: Jumlah tanggungan anggota keluarga nasabah.
    \item \texttt{Education\_Level}: Tingkat pendidikan terakhir yang ditempuh nasabah.
    \item \texttt{Marital\_Status}: Status pernikahan nasabah (contoh: \textit{Married, Single, Divorced, Unknown}).
    \item \texttt{Income\_Category}: Rentang estimasi pendapatan tahunan nasabah.
    \item \texttt{Card\_Category}: Jenis tingkatan kartu kredit yang dimiliki (Blue, Silver, Gold, Platinum).
    \item \texttt{Months\_on\_book}: Durasi atau lama waktu (dalam hitungan bulan) nasabah telah berlangganan layanan di bank.
    \item \texttt{Credit\_Limit}: Batas nominal maksimal kredit yang diberikan kepada kartu nasabah.
\end{itemize}


% ==========================================
% DATA PREPROCESSING
% ==========================================
\subsection{Data Preprocessing}

\begin{itemize}
    \item Dataset: https://www.kaggle.com/datasets/amitabhajoy/bengaluru-house-price-data
    \item Silahkan membaca dataset diatas menggunakan dataframe!
    \item Buatlah visualisasi histogram untuk beberapa variabel numerik, lalu jelaskan pola distribusi yang terlihat!
    \item Hitung dan visualisasikan korelasi antar variabel, kemudian jelaskan hubungan yang ditemukan!
    \item Lakukan exploratory data analysis (EDA) dengan membuat descriptive statistics serta melakukan pengecekan missing value dan duplikasi. Jelaskan tindakan apa yang perlu dilakukan jika keduanya ditemukan!
    \item Lakukan preprocessing lanjutan berupa encoding variabel kategorikal (one-hot \& label encoding), melakukan train–test split (80:20), serta menerapkan standardization atau normalization pada fitur numerik. Jelaskan alasan penggunaan masing-masing teknik!
\end{itemize}

\subsubsection{1. Impor Pustaka dan Membaca Dataset}
\begin{figure}[H]
    \centering
    \includegraphics[width=0.8\textwidth]{53.png}
    \caption{Source Code Impor Pustaka dan Load Dataset}
    \label{fig:53}
\end{figure}

\begin{figure}[H]
    \centering
    \includegraphics[width=0.8\textwidth]{54.png}
    \caption{Output df\_house.head() dan df\_house.tail()}
    \label{fig:54}
\end{figure}
Penjelasan Gambar :
\begin{itemize}
    \item \textbf{Input:} Memuat pustaka yang dibutuhkan dan mengunduh dataset \textit{Bengaluru House Price} dari Google Drive menggunakan file ID yang telah diperbarui. Dataset dibaca ke memori menggunakan \texttt{pd.read\_csv()}.
    \item \textbf{Output:} Menampilkan cuplikan 5 baris teratas dan terbawah, memperlihatkan rincian atribut struktur rumah seperti \texttt{area\_type}, \texttt{location}, \texttt{size}, \texttt{total\_sqft}, \texttt{bath}, \texttt{balcony}, dan label \texttt{price}.
\end{itemize}

\subsubsection{2. Pengecekan Descriptive Statistics, Missing Value, dan Duplikasi}
\begin{figure}[H]
    \centering
    \includegraphics[width=0.7\textwidth]{55.png}
    \caption{Source Code dan Output Pengecekan Data}
    \label{fig:55}
\end{figure}
Penjelasan Gambar :
\begin{itemize}
    \item \textbf{Input:} Menjalankan \texttt{describe()} untuk melihat sebaran statistik, serta \texttt{isnull().sum()} dan \texttt{duplicated().sum()} untuk mendeteksi anomali \textit{missing value} maupun duplikasi baris.
    \item \textbf{Output:} Terlihat tabel statistik univariat dasar. Ditemukan adanya 529 baris data duplikat dan \textit{missing value} yang bervariasi pada beberapa kolom, seperti \texttt{society} (5502 kosong), \texttt{balcony} (609 kosong), dan \texttt{bath} (73 kosong).
    \item \textbf{Tindakan yang Perlu Dilakukan:} 
    \begin{itemize}
        \item \textbf{Duplikasi:} Wajib dibersihkan menggunakan \texttt{drop\_duplicates()} agar model tidak mengalami bias (\textit{overfitting}) karena terlalu sering melihat data repetitif.
        \item \textbf{Missing Value:} Baris dengan nilai kosong dapat langsung dihapus (\texttt{dropna()}) jika persentasenya kecil. Alternatif lain adalah dengan mengisi kekosongan tersebut (\textit{imputasi}) menggunakan nilai \textit{median} untuk numerik atau \textit{modus} untuk kategorikal.
    \end{itemize}
\end{itemize}

\subsubsection{3. Visualisasi Histogram}
\begin{figure}[H]
    \centering
    \includegraphics[width=0.7\textwidth]{56.png}
    \caption{Source Code Visualisasi Histogram}
    \label{fig:56}
\end{figure}

\begin{figure}[H]
    \centering
    \includegraphics[width=0.8\textwidth]{57.png}
    \caption{Output Visualisasi Histogram}
    \label{fig:57}
\end{figure}
Penjelasan Gambar :
\begin{itemize}
    \item \textbf{Input:} Melakukan pembersihan data secara sementara (\textit{drop duplicates} dan \textit{drop missing values}), kemudian membuat plot histogram berukuran 12x8 untuk variabel \texttt{bath}, \texttt{balcony}, dan \texttt{price}.
    \item \textbf{Output \& Pola Distribusi:} Variabel \texttt{price} dan \texttt{bath} memperlihatkan pola distribusi \textit{Right-Skewed} (menjulur ke kanan). Hal ini berarti sebagian besar titik data properti berkumpul pada rentang harga menengah-bawah dengan fasilitas standar (sisi kiri grafik), sedangkan segelintir rumah mewah ekstrim membentuk ekor panjang ke kanan.
\end{itemize}

\subsubsection{4. Analisis Korelasi Antar Variabel}
\begin{figure}[H]
    \centering
    \includegraphics[width=0.7\textwidth]{58.png}
    \caption{Source Code dan Output Analisis Korelasi}
    \label{fig:58}
\end{figure}
Penjelasan Gambar :
\begin{itemize}
    \item \textbf{Input:} Menghitung matriks korelasi Pearson dengan \texttt{corr()} khusus untuk variabel numerikal yang difilter, kemudian dirender dalam bentuk peta panas (\textit{heatmap}) via \texttt{seaborn}.
    \item \textbf{Output \& Hubungan:} Visualisasi menunjukkan adanya korelasi positif moderat sebesar 0.46 antara jumlah kamar mandi (\texttt{bath}) dengan harga properti (\texttt{price}). Hubungan ini masuk akal; bertambahnya kamar mandi menandakan spesifikasi/luas rumah yang lebih besar, yang otomatis mengatrol nilai jual rumah tersebut.
\end{itemize}

\subsubsection{5. Preprocessing Lanjutan (Encoding, Splitting, dan Scaling)}
\begin{figure}[H]
    \centering
    \includegraphics[width=0.8\textwidth]{59.png}
    \caption{Source Code dan Output Preprocessing Lanjutan}
    \label{fig:59}
\end{figure}
Penjelasan Gambar :
\begin{itemize}
    \item \textbf{Input:} Menerapkan prapemrosesan tahap akhir yang mencakup: \textit{Label Encoding} pada kolom \texttt{area\_type}, \textit{One-Hot Encoding} pada \texttt{availability}, pemecahan himpunan data (\textit{Train-Test Split}), dan penstandaran nilai (\textit{Standardization}) menggunakan \texttt{StandardScaler()}.
    \item \textbf{Alasan Penggunaan Masing-masing Teknik:}
    \begin{itemize}
        \item \textbf{Label Encoding:} Digunakan untuk mengonversi kategori nilai string menjadi skema penomoran (0, 1, 2). Komputer memerlukan format numerik agar bisa menghitung matriks algoritma.
        \item \textbf{One-Hot Encoding:} Penting untuk diterapkan pada kategori nominal (kategori tanpa tingkatan superioritas). Memecahnya menjadi kolom biner mencegah algoritma salah tangkap dalam menganggap label "3" lebih tinggi bobotnya daripada "1".
        \item \textbf{Train-Test Split (80:20):} Memisahkan sebagian besar dataset (80\%) sebagai bahan ajar untuk model, dan menyisakan data baru (20\%) sebagai arena pengujian. Teknik ini vital guna mengevaluasi kualitas tebakan model tanpa terjebak \textit{overfitting}.
        \item \textbf{Standardization:} Menyamakan jangkauan ekstrem pada fitur. Teknik ini meratakan perbedaan besar (misal rasio harga jutaan dan jumlah kamar mandi satuan) agar fitur bernominal masif tidak merusak atau mendominasi kalkulasi jarak bobot pada saat pembelajaran mesin berlangsung.
    \end{itemize}
\end{itemize}

\chapter{Kesimpulan}
\label{chap:kesimpulan}
Berdasarkan praktikum \textit{Exploratory Data Analysis} (EDA) dan \textit{Data Preprocessing} yang telah dilaksanakan, diperoleh beberapa simpulan penting sebagai berikut:
\begin{enumerate}
    \item \textbf{Efektivitas EDA sebagai Diagnosa Data:} Melalui penerapan statistik deskriptif dan visualisasi (seperti histogram, \textit{box plot}, \textit{scatter plot}, dan \textit{heatmap}), karakteristik distribusi data, skewness, anomali \textit{missing value}, dan korelasi antar variabel dapat diidentifikasi secara tepat sebelum pemodelan dilakukan.
    
    \item \textbf{Peran Vital Data Preprocessing:} Penanganan anomali data melalui pembersihan duplikasi baris dan imputasi nilai kosong (\textit{missing values}) terbukti efektif menjaga integritas dataset tanpa kehilangan informasi penting secara berlebihan.
    
    \item \textbf{Optimalisasi Representasi Data via Encoding:} Transformasi atribut kategorikal menggunakan \textit{One-Hot Encoding} dan \textit{Label Encoding} berhasil menjembatani kebutuhan komputasi algoritma yang berbasis matriks numerik, sekaligus mencegah bias interpretasi tingkatan pada data nominal.
    
    \item \textbf{Stabilitas Skala dan Evaluasi Objektif:} Penerapan \textit{feature scaling} (Standardisasi dan Normalisasi) mampu menyetarakan rentang fitur numerik sehingga model tidak didominasi oleh variabel dengan nominal ekstrem. Selain itu, pembagian \textit{train-test split} (80:20) menjamin evaluasi performa model yang tidak bias terhadap data baru.
    
    \item \textbf{Pemanfaatan Ekosistem Python:} Integrasi pustaka Pandas, NumPy, Matplotlib, Seaborn, dan Scikit-Learn mempermudah eksekusi manipulasi data tabular dan visualisasi analitik secara terstruktur dan efisien.
\end{enumerate}


\chapter{Daftar Pustaka}
\begin{enumerate}
    \item \label{modulPraktikum} Dr. Eng. Ir. Ganjar Alfian, S.T., M. Eng., dan Tim Asisten. \textit{Modul Praktikum Penambangan Data: Exploratory Data Analysis (EDA) dan Data Preprocessing}. Yogyakarta: Departemen Teknik Elektro dan Informatika, Sekolah Vokasi, Universitas Gadjah Mada, 2026.
    
    \item \label{tukey} J. W. Tukey. \textit{Exploratory Data Analysis}. Reading, MA: Addison-Wesley, 1977.
    
    \item \label{mckinney} W. McKinney. \textit{Python for Data Analysis: Data Wrangling with Pandas, NumPy, and Jupyter}, 3rd ed. Sebastopol, CA: O'Reilly Media, 2022.
    
    \item \label{kaggleCredit} S. Goyal. \textit{Credit Card Customers Dataset}. Kaggle, 2020. [Online]. Tersedia: \url{https://www.kaggle.com/datasets/sakshigoyal7/credit-card-customers}.
    
    \item \label{kaggleBengaluru} A. Joy. \textit{Bengaluru House Price Data}. Kaggle, 2018. [Online]. Tersedia: \url{https://www.kaggle.com/datasets/amitabhajoy/bengaluru-house-price-data}.
    
    \item \label{colabLaporan} H. C. A. Paliling. \textit{S5-PPD-542159 (Source Code Laporan Praktikum)}. Google Colaboratory, 2026. [Online]. Tersedia: \url{https://colab.research.google.com/drive/1bvBJJTwzyzKHsH95cHhvi076hqc8DzNy?usp=sharing}.
\end{enumerate}

\end{document}